\chapter{Funções utilitárias de avaliação}\label{ap:evalutils}

O módulo \texttt{eval\_utils.py} concentra as funções reutilizadas por todos os experimentos do Capítulo~\ref{cap:investigacao}. As funções principais são reproduzidas abaixo.

\begin{lstlisting}[caption={Intervalo de confiança \textit{bootstrap} para ROC-AUC.},label={lst:bootstrap}]
def bootstrap_auc_ci(y_true, y_score, n_boot=1000,
                     alpha=0.05, seed=42):
    rng = np.random.default_rng(seed)
    n = len(y_true)
    aucs = np.empty(n_boot)
    for i in range(n_boot):
        idx = rng.integers(0, n, size=n)
        if len(np.unique(y_true[idx])) < 2:
            aucs[i] = np.nan
            continue
        aucs[i] = roc_auc_score(y_true[idx], y_score[idx])
    aucs = aucs[~np.isnan(aucs)]
    auc_point = roc_auc_score(y_true, y_score)
    lower = float(np.quantile(aucs, alpha / 2))
    upper = float(np.quantile(aucs, 1 - alpha / 2))
    return float(auc_point), lower, upper
\end{lstlisting}

\begin{lstlisting}[caption={Split temporal walk-forward 70/15/15.},label={lst:walkforward}]
def walk_forward_split(df, train_frac=0.70,
                       val_frac=0.15):
    n = len(df)
    n_train = int(n * train_frac)
    n_val = int(n * val_frac)
    train = df.iloc[:n_train]
    val = df.iloc[n_train:n_train + n_val]
    test = df.iloc[n_train + n_val:]
    return train, val, test
\end{lstlisting}

\chapter{Arquitetura TCN}\label{ap:tcn}

A TCN (\textit{Temporal Convolutional Network}) utiliza blocos de convolução causal dilatada com conexões residuais. A configuração TCN~[32,32] contém duas camadas com 32 filtros cada, \textit{kernel size} 3 e dilatações $[1, 2]$.

\begin{lstlisting}[caption={Bloco convolucional causal da TCN.},label={lst:tcn}]
class CausalConv1dBlock(nn.Module):
    def __init__(self, in_ch, out_ch, kernel_size,
                 dilation, dropout=0.2):
        super().__init__()
        padding = (kernel_size - 1) * dilation
        self.conv1 = nn.Conv1d(in_ch, out_ch,
            kernel_size, padding=padding,
            dilation=dilation)
        self.conv2 = nn.Conv1d(out_ch, out_ch,
            kernel_size, padding=padding,
            dilation=dilation)
        self.chomp = padding
        self.relu = nn.ReLU()
        self.dropout = nn.Dropout(dropout)
        self.downsample = (nn.Conv1d(in_ch, out_ch, 1)
                          if in_ch != out_ch else None)

    def forward(self, x):
        out = self.dropout(self.relu(
            self.conv1(x)[:, :, :-self.chomp]))
        out = self.dropout(self.relu(
            self.conv2(out)[:, :, :-self.chomp]))
        res = (self.downsample(x)
               if self.downsample else x)
        return self.relu(out + res)
\end{lstlisting}

\chapter{Exemplos de artigos e classificação de sentimento}\label{ap:artigos}

Este apêndice ilustra, com artigos reais do corpus InfoMoney, como o FinBERT-PT-BR \cite{souza2023} converte o texto de uma notícia nas \textit{features} de sentimento usadas pelos modelos. Para cada artigo, o modelo produz três \textit{logits} --- pontuações brutas não normalizadas associadas às classes POSITIVO, NEGATIVO e NEUTRO. Aplicando a função \textit{softmax} aos \textit{logits} obtêm-se probabilidades que somam~1; a classe atribuída é a de maior probabilidade (\textit{argmax}). A coluna \texttt{mean\_logit\_pos/neg/neu} do \textit{dataset} (Seção~\ref{cap:sentimento}) é a média diária desses \textit{logits} sobre todos os artigos do dia, e \texttt{mean\_sentiment} é o código numérico da classe predominante.

A Tabela~\ref{tab:exemplos_artigos} reproduz seis manchetes reais do corpus --- duas por classe --- com seus \textit{logits} e probabilidades efetivamente computados pelo modelo. As entradas foram extraídas dos arquivos \texttt{petr4\_noticias.json}, \texttt{vale3\_noticias.json} e \texttt{itub4\_noticias.json} e processadas exatamente como no \textit{pipeline} de extração (entrada \texttt{título}, truncamento em 512 \textit{tokens}).

\begin{table}[htb]
\centering
\caption{Manchetes reais do corpus InfoMoney e respectiva classificação de sentimento pelo FinBERT-PT-BR. \textit{Logits} na ordem (POS, NEG, NEU); probabilidades obtidas via \textit{softmax}. A classe atribuída é o \textit{argmax}.}\label{tab:exemplos_artigos}
\footnotesize
\begin{tabularx}{\textwidth}{Xccc}
\toprule
\textbf{Manchete} & \textbf{\textit{Logits} (POS,\,NEG,\,NEU)} & \textbf{Prob.\ (POS,\,NEG,\,NEU)} & \textbf{Classe} \\
\midrule
``Hypera projeta aceleração das vendas e tem semaglutida como principal aposta em 2026'' & $(1{,}33;\ -1{,}85;\ -1{,}36)$ & $(0{,}90;\ 0{,}04;\ 0{,}06)$ & POSITIVO \\
\addlinespace
``Com maior peso de estatais na AL, Brasil abre oportunidade assimétrica na B3, diz BBI'' & $(0{,}88;\ -2{,}07;\ -0{,}88)$ & $(0{,}82;\ 0{,}04;\ 0{,}14)$ & POSITIVO \\
\addlinespace
``Ibovespa fecha abaixo de 178 mil pontos sem alívio em preocupações com guerra no Irã'' & $(-1{,}73;\ 1{,}45;\ -1{,}80)$ & $(0{,}04;\ 0{,}93;\ 0{,}04)$ & NEGATIVO \\
\addlinespace
``Ibovespa fecha com nova queda e emenda terceira semana seguida no negativo'' & $(-1{,}59;\ 1{,}41;\ -1{,}82)$ & $(0{,}05;\ 0{,}92;\ 0{,}04)$ & NEGATIVO \\
\addlinespace
``Mini-índice (WINJ26): petróleo, PIB dos EUA e inflação guiam traders hoje'' & $(-1{,}45;\ -1{,}30;\ 0{,}91)$ & $(0{,}08;\ 0{,}09;\ 0{,}83)$ & NEUTRO \\
\addlinespace
``Mini-índice (WINJ26): guerra, petróleo e agenda da semana guiam traders'' & $(-1{,}69;\ -0{,}99;\ 0{,}80)$ & $(0{,}07;\ 0{,}13;\ 0{,}80)$ & NEUTRO \\
\bottomrule
\end{tabularx}
\end{table}

Os exemplos evidenciam o comportamento esperado: manchetes com projeções de crescimento ou oportunidade recebem alta probabilidade POSITIVA; quedas e quadros de aversão a risco, alta probabilidade NEGATIVA; e manchetes meramente descritivas de agenda, alta probabilidade NEUTRA. Vale notar, contudo, que a confiança do modelo é sensível à formulação da manchete --- um mesmo evento econômico pode receber rótulos distintos conforme o enquadramento jornalístico. Essa variância na rotulagem de manchetes curtas é coerente com o ruído de sinal discutido no Capítulo~\ref{cap:investigacao} e ajuda a explicar por que a agregação diária dos \textit{logits} agrega, na prática, pouca informação preditiva sobre a direção dos preços.
